\section{模型框架}

本节构建一个融合新结构经济学视角的异质企业行业均衡模型，用以分析药品上市许可持有人制度（`MAH`）如何通过改变研发所有权与生产能力之间的制度边界，影响企业分化、组织重构与原研药创新。模型的基本思想是：发展阶段决定制度需求，政府通过制度供给改变企业面临的制度摩擦，市场则通过企业进入、退出、分工和组织状态转移完成资源重新配置，最终影响创新与福利。

\subsection{发展阶段、制度供给与市场反应}

本文以新结构经济学为上层理论视角。设 $\xi \in \Xi \subset \mathbb{R}_{+}$ 表示医药产业所处的发展阶段。较高的 $\xi$ 表示产业结构更接近原研药创新导向，研发、项目推进和专业化分工的需求更强。政府在给定 $\xi$ 的条件下选择制度安排
\[
M \in \{0,1\},
\]
其中 $M=0$ 表示改革前制度，$M=1$ 表示 `MAH` 制度。与将政策视为纯粹外生冲击不同，本文强调制度安排是否适宜于发展阶段需求；但在操作上，本文仍将 $M$ 视为给定制度环境，而不刻画政府与企业之间的动态博弈。

行业中存在三类活跃企业和一类潜在进入者：`A` 型企业同时从事研发与生产，`B` 型企业专门从事研发，`C` 型企业专门从事生产，`E` 为潜在进入者。每个活跃企业拥有能力状态 $z \in Z \subset \mathbb{R}_{+}$，该状态综合反映企业的研发能力、项目推进能力、组织协同能力以及商业化和重组能力。

为了将“制度是否适宜发展阶段”写入模型，设制度适配函数为
\[
\Lambda(z,\xi)=z^{\alpha}\xi^{\nu}, \qquad \alpha>0,\ \nu>0.
\]
其中，$\Lambda(z,\xi)$ 表示能力为 $z$ 的企业在发展阶段 $\xi$ 下利用制度改革的能力。由定义可得
\[
\frac{\partial \Lambda}{\partial z}>0, \qquad \frac{\partial \Lambda}{\partial \xi}>0.
\]
因此，企业能力越强、发展阶段越高，制度改革越可能转化为真实的创新和组织重构收益。

\subsection{制度摩擦与组织重构门槛}

`MAH` 的作用在模型中主要通过两个渠道体现。第一，降低研发型企业将原研药 idea 推向商业化的制度摩擦。设该摩擦为
\[
\tau(z,\xi,M)=\bar{\tau}-M\Delta_{\tau}\Lambda(z,\xi), \qquad \bar{\tau}>0,\ \Delta_{\tau}>0.
\]
则有
\[
\frac{\partial \tau}{\partial M}=-\Delta_{\tau}\Lambda(z,\xi)<0,
\qquad
\frac{\partial^{2}\tau}{\partial M\partial \xi}
=-\Delta_{\tau}\frac{\partial \Lambda}{\partial \xi}<0.
\]
这表明，`MAH` 能够降低商业化摩擦，且这种降低在发展阶段更高时更为显著。

第二，`MAH` 降低企业在不同组织形式之间切换的重组成本。设企业从组织状态 $s$ 转向 $j$ 的重组成本为
\[
\kappa_{sj}(z,\xi,M)=\bar{\kappa}_{sj}-M\Delta_{sj}\Lambda(z,\xi), \qquad \Delta_{sj}>0.
\]
于是
\[
\frac{\partial \kappa_{sj}}{\partial M}<0,
\qquad
\frac{\partial^{2}\kappa_{sj}}{\partial M\partial \xi}<0.
\]
这意味着，制度改革不仅降低企业组织重构的门槛，而且这种重构激励在制度与发展阶段更匹配时更强。

\subsection{企业利润函数与原研药创新决策}

本文中的 innovation object 统一指向原研药创新。`A` 型企业同时拥有研发和生产能力，给定工资 $w$ 和创新强度 $x_A \ge 0$，其流量收益为
\[
\pi_A(z,\xi;w)
=
\bar{\pi}_A(z,\xi;w)-f_A+x_A\Omega_A(z,\xi)-\frac{\chi_A}{\eta}x_A^{\eta},
\qquad \eta>1.
\]
由一阶条件
\[
\Omega_A(z,\xi)-\chi_A x_A^{\eta-1}=0
\]
可得最优创新强度
\[
x_A^*(z,\xi)
=
\left(\frac{\Omega_A(z,\xi)}{\chi_A}\right)^{\frac{1}{\eta-1}},
\]
并得到最优流量收益
\[
\Pi_A(z,\xi;w)
=
\bar{\pi}_A(z,\xi;w)-f_A
+ \frac{\eta-1}{\eta}\chi_A^{-\frac{1}{\eta-1}}
\Omega_A(z,\xi)^{\frac{\eta}{\eta-1}}.
\]

`B` 型企业不进行内部生产，其核心决策是选择 idea generation 强度 $x_B \ge 0$ 并决定 idea 如何进入 implementation/commercialization 阶段。设每个 idea 的净价值为 $\Psi_B(z,\xi;p_m,M)$，则其流量收益为
\[
\pi_B(z,\xi;p_m,M)
=
-f_B+x_B\Psi_B(z,\xi;p_m,M)-\frac{\chi_B}{\eta}x_B^{\eta}.
\]
由一阶条件
\[
\Psi_B(z,\xi;p_m,M)-\chi_B x_B^{\eta-1}=0
\]
可得
\[
x_B^*(z,\xi;p_m,M)
=
\left(\frac{\Psi_B(z,\xi;p_m,M)}{\chi_B}\right)^{\frac{1}{\eta-1}},
\]
从而
\[
\Pi_B(z,\xi;p_m,M)
=
-f_B+\frac{\eta-1}{\eta}\chi_B^{-\frac{1}{\eta-1}}
\Psi_B(z,\xi;p_m,M)^{\frac{\eta}{\eta-1}}.
\]

`C` 型企业提供委托生产服务。若服务供给量为 $m \ge 0$，则
\[
\pi_C(z,\xi;p_m,w)=p_m m-c(m,z,\xi;w)-f_C,
\]
其中 $c_m>0$、$c_{mm}>0$、$c_z<0$。相应最优利润为
\[
\Pi_C(z,\xi;p_m,w)
=
\max_{m\ge 0}\left\{p_m m-c(m,z,\xi;w)-f_C\right\}.
\]

\subsection{Invention 与 Implementation 的分离}

为了刻画 `MAH` 如何促进原研药创新，本文将 innovation 分为 invention 与 implementation 两个阶段。设 `B` 型企业生成一个 idea 后，面临三种选择：保持研发主体身份并委托 `C` 型企业生产；支付组织重构成本转为 `A` 型企业后内部实施；或者不实施该 idea。

若选择委托生产，则每个 idea 的净价值为
\[
\Psi_B^O(z,\xi;p_m,M)
=
\max_{i_O \in [0,1]}
\left\{
i_O R_O(z,\xi)-\frac{\kappa_O}{\theta}i_O^\theta-p_m-\tau(z,\xi,M)
\right\},
\qquad \theta>1.
\]
由一阶条件
\[
R_O(z,\xi)-\kappa_O i_O^{\theta-1}=0
\]
可得
\[
i_O^*(z,\xi)=\left(\frac{R_O(z,\xi)}{\kappa_O}\right)^{\frac{1}{\theta-1}}.
\]
将其代回得到
\[
\Psi_B^O(z,\xi;p_m,M)
=
\frac{\theta-1}{\theta}\kappa_O^{-\frac{1}{\theta-1}}
R_O(z,\xi)^{\frac{\theta}{\theta-1}}
-p_m-\tau(z,\xi,M).
\]
因此，
\[
\frac{\partial \Psi_B^O}{\partial M}
=-\frac{\partial \tau}{\partial M}
=\Delta_{\tau}\Lambda(z,\xi)>0,
\]
且
\[
\frac{\partial^2 \Psi_B^O}{\partial M \partial \xi}
=\Delta_{\tau}\frac{\partial \Lambda}{\partial \xi}>0.
\]
这说明 `MAH` 对研发型企业外包商业化价值的提升，在更高发展阶段上更强。

若 `B` 型企业转为 `A` 型企业后内部实施，则每个 idea 的净价值为
\[
\Psi_B^I(z,\xi;M)
=
\max_{i_I \in [0,1]}
\left\{
i_I R_I(z,\xi)-\frac{\kappa_I}{\theta}i_I^\theta-\kappa_{BA}(z,\xi,M)
\right\}.
\]
类似地可得最优内部实施强度
\[
i_I^*(z,\xi)=\left(\frac{R_I(z,\xi)}{\kappa_I}\right)^{\frac{1}{\theta-1}}.
\]
因此，`B` 型企业每个 idea 的总价值为
\[
\Psi_B(z,\xi;p_m,M)
=
\max\left\{
\Psi_B^O(z,\xi;p_m,M),\,
\Psi_B^I(z,\xi;M),\,
0
\right\}.
\]
这一对象同时承载三层含义：其一，创新并不自动从 invention 变成产出；其二，`MAH` 通过降低商业化制度摩擦提高研发型企业的 implementation 价值；其三，这种作用依赖于制度与企业能力结构、发展阶段之间的匹配程度。

\subsection{组织状态转移}

企业当前组织状态为 $s \in \{A,B,C\}$，下一期可选状态集合记为 $\mathcal{J}(s)$。若企业当前类型为 $s$，选择下一期状态为 $j \in \mathcal{J}(s)$，则系统性价值为
\[
\bar V_{sj}(z,\xi;M)
=
\Pi_j(z,\xi;M)-\kappa_{sj}(z,\xi,M)
+\beta E[V_j(z',\xi;M)\mid z,s,j].
\]
若企业选择时还受到 i.i.d. type-I extreme value shocks $\varepsilon_{sj}$ 影响，则 Bellman 方程为
\[
V_s(z,\xi;M)
=
E_{\varepsilon}
\left[
\max_{j\in \mathcal{J}(s)}
\left\{
\bar V_{sj}(z,\xi;M)+\varepsilon_{sj}
\right\}
\right].
\]
根据标准 logit 公式，组织转移概率为
\[
P_{sj}(z,\xi;M)
=
\frac{\exp(\bar V_{sj}(z,\xi;M)/\sigma)}
{\sum_{\ell \in \mathcal{J}(s)}\exp(\bar V_{s\ell}(z,\xi;M)/\sigma)}.
\]
由于 $\partial \kappa_{sj}/\partial M<0$ 且 $\partial^2 \kappa_{sj}/\partial M \partial \xi<0$，对于关键转移方向如 `A \to B`、`A \to C`、`B \to A` 和 `C \to A`，一般有
\[
\frac{\partial P_{sj}}{\partial M}>0,
\qquad
\frac{\partial^2 P_{sj}}{\partial M \partial \xi}>0.
\]
这意味着，在更高发展阶段上，`MAH` 对企业分化和组织重构的诱导作用更强。

\subsection{进入、市场清算与稳态分布}

潜在进入者初始能力服从分布 $G(z_0)$，其自由进入条件为
\[
\int V_B(z_0,\xi;M)\,dG(z_0)\le c_{EB},
\qquad
\int V_C(z_0,\xi;M)\,dG(z_0)\le c_{EC}.
\]
如果存在正进入，则相应条件取等号。由于 `MAH` 通过提升 $\Psi_B$ 进而提升 $V_B$，故研发型进入对制度改革和发展阶段均具有系统性响应。

`B` 型企业的委托生产需求与 `C` 型企业的服务供给分别为
\[
D_m(p_m,\xi;M)
=
\int x_B^*(z,\xi;p_m,M)d_O(z,\xi;p_m,M)\mu_B(dz),
\]
\[
S_m(p_m,\xi;w)
=
\int m^*(z,\xi;p_m,w)\mu_C(dz),
\]
市场清算条件为
\[
D_m(p_m,\xi;M)=S_m(p_m,\xi;w).
\]

记稳态分布为 $(\mu_A,\mu_B,\mu_C)$，其满足由进入、退出、状态转移和能力演化共同决定的 Kolmogorov 型平衡条件。

\subsection{原研药创新、福利与适宜制度}

本文将创新数量、创新质量和创新结构统一理解为原研药创新对象。行业原研药创新产出为
\[
\mathcal{I}(M,\xi)
=
\int \Gamma_A(z,\xi)x_A^*(z,\xi)\mu_A(dz)
+\int \Gamma_B(z,\xi)x_B^*(z,\xi;p_m,M)i_B^*(z,\xi;p_m,M)\mu_B(dz).
\]
福利记为
\[
\mathcal{W}(M,\xi).
\]
在该框架下，制度改革的净效应同时受到四类力量影响：原研药创新增益、企业分化与资源重配置收益、传统生产企业利润损失以及组织调整成本。因此，创新增加并不自动意味着福利增加。

为了体现新结构经济学中的“有为政府提供适宜制度”，设政府在观察发展阶段 $\xi$ 后，通过最大化净福利来选择制度：
\[
\max_{M \in \{0,1\}} \mathcal{W}(M,\xi).
\]
定义制度改革的净收益为
\[
\Delta \mathcal{W}(\xi)=\mathcal{W}(1,\xi)-\mathcal{W}(0,\xi).
\]
若存在阈值 $\xi^*$，使得
\[
\Delta \mathcal{W}(\xi)<0,\quad \xi<\xi^*,
\]
\[
\Delta \mathcal{W}(\xi)\ge 0,\quad \xi\ge \xi^*,
\]
则说明 `MAH` 只有在产业发展阶段达到一定水平后才成为适宜制度。换言之，制度改革对创新和产业升级的作用，本质上取决于制度供给与发展阶段、企业能力结构之间的匹配程度。

\subsection{主要命题}

基于上述框架，本文得到以下几个可检验命题。第一，`MAH` 对研发型企业商业化价值的提升，在更高发展阶段和更高能力企业上更强。第二，`MAH` 通过降低组织重构门槛诱导企业分化与状态转移。第三，原研药创新增加来自 `B` 型进入、implementation 改善、既有企业创新激励变化以及组织结构重组。第四，制度促进原研药创新的作用具有显著的发展阶段异质性。第五，净福利为正是条件性结论：只有当创新增益与资源重配置收益超过传统生产利润损失和组织调整成本时，`MAH` 才带来正的净福利。

\subsection{与实证的对应}

该模型与实证的对应关系十分直接。企业分化与组织状态转移可由 `A/B/C` 的转移矩阵来刻画；原研药创新数量、质量和结构分别对应模型中的 $\mathcal{I}(M,\xi)$ 及其构成；`C` 型企业利润变化对应传统生产企业利润受损；而地区产业基础、医药集聚度、企业初始能力结构等变量则可以用来刻画发展阶段 $\xi$ 的异质性，从而检验“适宜制度安排”这一新结构经济学命题。
